\documentclass[12pt]{article}

\usepackage[T1]{fontenc}
\usepackage[utf8]{inputenc}
\usepackage{lmodern}
\usepackage[margin=1in]{geometry}
\usepackage{microtype}
\usepackage{amsmath,amssymb,amsthm,mathtools}
\usepackage{array,booktabs,enumitem}
\usepackage[colorlinks=true,linkcolor=blue,citecolor=blue,urlcolor=blue]{hyperref}
\usepackage[capitalize,nameinlink,noabbrev]{cleveref}

\newtheorem{theorem}{Theorem}[section]
\newtheorem{proposition}[theorem]{Proposition}
\newtheorem{lemma}[theorem]{Lemma}
\newtheorem{corollary}[theorem]{Corollary}
\newtheorem{condition}[theorem]{Condition}
\newtheorem{assumption}[theorem]{Assumption}
\theoremstyle{definition}
\newtheorem{definition}[theorem]{Definition}
\newtheorem{remark}[theorem]{Remark}
\newtheorem{example}[theorem]{Example}

\newcommand{\eps}{\varepsilon}
\newcommand{\R}{\mathbb{R}}
\newcommand{\N}{\mathbb{N}}
\newcommand{\E}{\mathbb{E}}
\newcommand{\Pb}{\mathbb{P}}
\newcommand{\1}{\mathbf{1}}
\newcommand{\cA}{\mathcal{A}}
\newcommand{\cB}{\mathcal{B}}
\newcommand{\cE}{\mathcal{E}}
\newcommand{\cF}{\mathcal{F}}
\newcommand{\cM}{\mathcal{M}}
\newcommand{\cN}{\mathcal{N}}
\newcommand{\cP}{\mathcal{P}}
\newcommand{\cQ}{\mathcal{Q}}
\newcommand{\cT}{\mathcal{T}}
\newcommand{\cS}{\mathcal{S}}
\newcommand{\cX}{\mathcal{X}}
\DeclareMathOperator{\conv}{conv}
\DeclareMathOperator{\supp}{supp}
\DeclareMathOperator{\spn}{sp}
\DeclareMathOperator{\Out}{Out}
\DeclareMathOperator{\BR}{BR}

\title{\textbf{Uniform Equilibrium from an Augmented Metastable Tree\\under Robust Nodewise Irreducibility}}
\author{Pedro Aldighieri}
\date{March 2026}

\begin{document}
\maketitle



\begin{abstract}
Neyman's problem asks whether every finite stochastic game admits a uniform
$\eps$-equilibrium. The full conjecture remains open. This paper proves the
\emph{downstream} part of the compact-tree route under explicit conditional input:
we assume that a finite universal $\eta$-tree representation is available together
with the first-exit controller realization attached to that tree, and we show that
this input yields a public finite-memory uniform $\eps$-equilibrium under a robust
nodewise irreducibility hypothesis.

Version~9 keeps the v8 repairs and fixes the three remaining bookkeeping
points from the conditional review: the gain/continuation coordinates are placed
in depth-dependent compact boxes, the finitely many prelude blocks before the
tail scale $M_*$ are peeled off into the $K_\eta/T$ remainder, and public
routing now specifies the terminal-label case explicitly. The bookkeeping kernel
with restart through $\rho_C$ is used only to define the stationary rate package
$\bar\mu_C^*=(\mu_C^*,\beta_C^*)$, and in v9 that restart is fixed at the
reference state $\xi_C^\circ$. The executed controller, by contrast, produces a
genuine first-exit renewal block of random duration $\tau\le L$, where $L$ is
only a truncation cap. The correct local identity is therefore the
\emph{renewal-reward relation}
\[
g_{i,C}^*
=
\frac{\E[\widetilde U_{i,C,L}]}{\E[\tau_{C,L}]}
\]
up to the controller realization error, and not the false scheduled-length
normalization $\E[\widetilde U_{i,C,L}/L]$.

To keep the convex occupation package and the genuine executed dynamics separate,
the paper distinguishes two kernels. The \emph{bookkeeping kernel} lives on
$\Xi_C$ alone and folds exit mass back to $\Xi_C$ through the internal restart map
$\rho_C$; it is used only to define the convex occupation package
$\bar\mu_C^*=(\mu_C^*,\beta_C^*)$. The \emph{executed augmented kernel} lives on
$\Xi_C^{\mathrm{aug}}$ and has absorbing cemetery states; it is used for
condition~\textnormal{(a)}, for the bias bound, and for the blockwise Bellman
estimate.

At each node and for frozen continuation values, every player faces a finite
communicating average-reward MDP on the internal state space with terminal
continuation rewards at exits. The normalized bias is unique after fixing one
reference coordinate, local stationary Nash equilibria exist by Kakutani, and the
global frozen-continuation correspondence has a fixed point. Using the first-exit
controller realization supplied by the tree input, the realized profile satisfies a
regret bound of the form
\[
\eta + C_\eta\delta + \frac{K_\eta}{T}.
\]
Choosing the parameters in that order yields the theorem.
\end{abstract}


\tableofcontents

\section{Introduction}

\subsection{Neyman's problem}

Shapley introduced stochastic games in 1953 and proved existence of stationary
equilibria for the discounted evaluation \cite{Shapley1953}. The long-run average
evaluation is much subtler. For zero-sum stochastic games, Mertens and Neyman
proved existence of the uniform value \cite{MertensNeyman1981}. For two-player
non-zero-sum stochastic games, Vieille established the existence of uniform
$\eps$-equilibria \cite{Vieille2000I,Vieille2000II}. Beyond two players the full
problem remains open; see Neyman's survey chapter, Solan's chapter on the
multi-player problem, Vieille's survey, and the handbook treatment
\cite{Neyman2003,Solan2003,Vieille2002,MertensSorinZamir2003}. Positive results
are known in special subclasses, including three-player absorbing games
\cite{Solan1999}, correlated equilibrium constructions \cite{SolanVieille2002},
cyclic Markov equilibria phenomena \cite{FleschThuijsmanVrieze1997}, and extensions
to broader state-space settings \cite{Nowak2003}.

The central question is:
\begin{quote}
\textbf{Neyman's open problem.}
For every finite stochastic game and every $\eps>0$, does there exist a strategy
profile that is an $\eps$-equilibrium simultaneously for all sufficiently large
finite horizons?
\end{quote}

This paper does not claim to settle the conjecture in full generality. Its scope
is narrower and more precise: we start from the finite universal $\eta$-tree
representation produced by the route-D' reduction, together with its first-exit
controller realization, and show that once that input is available the remainder
of the route is correct under a robust nodewise irreducibility condition. The
point of the rewrite is not to alter the theorem's scope but to make every
definition used downstream mathematically correct and every stated proof complete.

\subsection{Main theorem}

We first record the standing scope. The compact asymptotic tree representation and
its first-exit controller realization are inputs to the present paper.

\begin{assumption}[Standing tree input]
\label{ass:tree-input}
For every $\eta\in(0,1)$, the game is endowed with a finite universal
$\eta$-tree representation as specified in
\Cref{ass:node-data,ass:controller-input}. All theorems in
Sections~\ref{sec:bookkeeping}--\ref{sec:proof-main} are proved relative to that
finite-dimensional input.
\end{assumption}

Under this standing input the paper proves the following theorem.

\begin{theorem}[Main theorem]
\label{thm:main}
Let
\[
G=(N,S,(A_i)_{i\in N},u,P)
\]
be a finite stochastic game, and assume \Cref{ass:tree-input}. Assume further
that condition~\textnormal{(a)} from \Cref{cond:a} holds for the augmented node
kernels of the universal $\eta$-trees. Then for every $\eps>0$ there exist a
public finite-memory strategy profile $\sigma^\eps$ and $T_0\in\N$ such that for
every $T\ge T_0$, every initial state $s\in S$, every player $i\in N$, and every
unilateral deviation $\tau_i$,
\[
\gamma_i^T(s,\sigma^\eps)
\ge
\gamma_i^T\bigl(s,(\tau_i,\sigma^\eps_{-i})\bigr)-\eps.
\]
In particular, every finite stochastic game equipped with a universal
$\eta$-tree representation and first-exit controller realization satisfying
condition~\textnormal{(a)} admits a public finite-memory uniform
$\eps$-equilibrium for every $\eps>0$.
\end{theorem}

\begin{remark}
The theorem is intentionally scoped to the downstream part of the compact-tree
route. The full Neyman conjecture remains open because the route still requires,
in the background, construction of the universal $\eta$-tree representation and
its realization map for an arbitrary game. The purpose of the present rewrite is
to show that the remaining argument from that representation to a uniform
equilibrium is mathematically correct once the earlier blockers and the
renewal-reward bookkeeping repair are in place.
\end{remark}


\subsection{The v9 repair}

Version~9 keeps the v8 frozen-continuation, local-Nash, bias-normalization,
augmented-kernel, routing, and renewal-reward repairs, and adds the final
bookkeeping fixes from the conditional review.

\begin{enumerate}[label=\textnormal{(\arabic*)},leftmargin=2.4em]
\item The bias normalization is linearized. Every local bias is normalized by the
fixed reference coordinate $h_{i,C}(\xi_C^\circ)=0$, not by the non-convex
condition $\min_\xi h_{i,C}(\xi)=0$.
\item Every irreducibility, hitting-time, and bias-span statement is formulated
for the genuine absorbing augmented kernel on
$\Xi_C^{\mathrm{aug}}=\Xi_C\sqcup\partial B_C$, not for the sub-stochastic
internal kernel.
\item A block at node $C$ terminates at its first cemetery visit. It emits at
most one boundary label, moves to the corresponding child node when the label
has a child successor, switches to the attached terminal continuation profile
when the label is terminal, and never regenerates through cemetery states.
\item The bookkeeping kernel and the executed first-exit dynamics are kept
separate. The restart kernel $\rho_C$ belongs only to the bookkeeping occupation
package, is never an executed transition after a cemetery hit, and is fixed by
the bookkeeping convention
\[
\rho_C(\cdot\mid b)=\delta_{\xi_C^\circ}
\qquad (b\in B_C).
\]
\item The local payoff identity is written in renewal-reward form. If a block is
launched with cap $L$, then its \emph{actual} duration is the random first-exit
time $\tau\le L$, and the correct target is
\[
g_{i,C}^*
=
\frac{\E[\widetilde U_{i,C,L}]}{\E[\tau_{C,L}]},
\]
not $\E[\widetilde U_{i,C,L}/L]$. The only possible extra bookkeeping term is the
restart-bias
\[
\sum_{\xi'}(R_C\beta_C^*)(\xi')h_{i,C}^*(\xi'),
\]
and in v9 it vanishes because every bookkeeping restart returns to the reference
state $\xi_C^\circ$ where $h_{i,C}^*(\xi_C^\circ)=0$. In the horizon-$T$
accounting, one sums the actual block durations and uses
$\sum_B |B| = T + O(1)$ rather than the scheduled caps.
\item The ambient compact boxes are depth-sensitive. For each node $C$, gains are
confined to $[0,d_\eta(C)]$ and boundary continuation values to
$[0,d_\eta(C)-1]^N$, where $d_\eta(C)$ is the remaining tree depth below $C$.
This is the correct box because each node contributes at most one unit of stage
payoff and continuation values propagate through only finitely many further
layers.
\item The finitely many prelude blocks before the frozen tail scale $M_*$ are not
covered by the tail block-regret estimate. Their total contribution is therefore
peeled off once and for all into the $K_\eta/T$ remainder term.
\item Public routing specifies the terminal-label case explicitly: if a cemetery
label $b$ leads to a terminal label rather than a child node, block scheduling
stops and play continues according to the terminal continuation profile attached
to $b$.
\end{enumerate}

\subsection{Roadmap}


Section~\ref{sec:model} recalls the game model and states the finite-dimensional
node data used by the compact-tree route. Section~\ref{sec:bookkeeping}
introduces the bookkeeping kernel and proves the graph theorem for the internal
occupation package. Section~\ref{sec:local} formulates the local game on the
genuine augmented state space and establishes local best-response and local
equilibrium existence. Section~\ref{sec:global} defines the global correspondence
with frozen continuation values and proves a Kakutani fixed point theorem.
Section~\ref{sec:realization} defines first-exit blocks and realizes the fixed
point by public finite-memory block controllers. Section~\ref{sec:proof-main}
proves the augmented-kernel bias bound, the deviation-value Lipschitz lemma, the
blockwise Bellman bound, and finally \Cref{thm:main}. The appendices record the
finite-MDP facts and the final aggregation bookkeeping.

\section{Model and standing asymptotic input}
\label{sec:model}

\subsection{Finite stochastic games}

A finite stochastic game is a tuple
\[
G=(N,S,(A_i)_{i\in N},u,P)
\]
where $N=\{1,\dots,n\}$ is the finite set of players, $S$ is the finite state
space, $A_i(s)$ is the finite action set of player $i$ at state $s$, the stage
payoff of player $i$ is a function
\[
u_i:S\times \prod_{j\in N}A_j(s)\to [0,1],
\]
and the transition kernel is
\[
P(\cdot\mid s,a)\in \Delta(S)
\qquad (s\in S,\ a\in \prod_{j\in N}A_j(s)).
\]
The public history at stage $t$ is
\[
h_t=(s_1,a_1,s_2,a_2,\dots,s_t).
\]
A behavioral strategy of player $i$ assigns to every public history $h_t$ a
probability distribution on $A_i(s_t)$.

For a horizon $T\ge 1$, an initial state $s$, and a profile $\sigma$, the average
payoff of player $i$ is
\[
\gamma_i^T(s,\sigma)
=
\E_{s,\sigma}\Bigl[\frac1T\sum_{t=1}^T u_i(s_t,a_t)\Bigr].
\]

\begin{definition}
A profile $\sigma$ is an \emph{$\eps$-equilibrium for horizon $T$} if for every
player $i$, every initial state $s$, and every unilateral deviation $\tau_i$,
\[
\gamma_i^T(s,\sigma)
\ge
\gamma_i^T\bigl(s,(\tau_i,\sigma_{-i})\bigr)-\eps.
\]
It is a \emph{uniform $\eps$-equilibrium} if there exists $T_0$ such that the same
inequality holds for every $T\ge T_0$.
\end{definition}

\subsection{Node data from the tree representation}

The present paper starts after the route-D' reduction has produced a finite rooted
universal $\eta$-tree. We record exactly the node data used later.

\begin{assumption}[Finite node data]
\label{ass:node-data}
Fix $\eta\in(0,1)$. The universal $\eta$-tree $\cT_\eta$ is finite. Every node
$C\in \cT_\eta$ carries the following data.
\begin{enumerate}[label=\textnormal{(D\arabic*)},leftmargin=2.8em]
\item A finite phase-state set
\[
\Xi_C=\bigsqcup_{k\in\mathbb Z/d_C\mathbb Z}\{k\}\times S_C^k
\]
and a distinguished reference state $\xi_C^\circ\in\Xi_C$.
\item For each player $i$ and phase-state $\xi\in\Xi_C$, a finite action set
$A_i(\xi)$.
\item A finite boundary set $B_C$ and a successor map $b\mapsto C[b]$ sending
boundary states to child nodes or terminal labels. If $C[b]$ is a terminal
label, the tree input also specifies a terminal continuation profile whose
payoff vector is denoted by $w(b)\in[0,1]^N$.
\item Internal transition coefficients
\[
P_C^0(\xi'\mid \xi,a)\in[0,1],
\qquad \xi,\xi'\in \Xi_C,
\ a\in \prod_{j\in N}A_j(\xi),
\]
and exit coefficients
\[
\lambda_C(b\mid \xi,a)\in[0,1],
\qquad b\in B_C,
\]
such that
\[
\sum_{\xi'\in\Xi_C}P_C^0(\xi'\mid \xi,a)+\sum_{b\in B_C}\lambda_C(b\mid\xi,a)=1.
\]
\item A bookkeeping restart kernel fixed by the convention
\[
\rho_C(\cdot\mid b)=\delta_{\xi_C^\circ}
\qquad (b\in B_C).
\]
This kernel is used only in the internal bookkeeping chain from
Section~\ref{sec:bookkeeping}; it is not part of the executed first-exit
dynamics. In particular every bookkeeping restart returns to the reference
state $\xi_C^\circ$.
\item A universal bias box constant $B_\eta^0<\infty$ such that every local
Bellman certificate used later can be normalized inside
$[-B_\eta^0,B_\eta^0]^{\Xi_C}$.
\end{enumerate}
\end{assumption}

\begin{remark}[Tree-depth bookkeeping]
\label{rem:tree-depth}
For each terminal label $\ell$, set $d_\eta(\ell):=1$. For a node $C\in\cT_\eta$,
define recursively
\[
d_\eta(C)
:=
\begin{cases}
1, & B_C=\varnothing,\\[0.3em]
1+\max_{b\in B_C} d_\eta(C[b]), & B_C\neq \varnothing,
\end{cases}
\]
where $d_\eta(C[b])=d_\eta(D)$ if $C[b]=D$ is a child node and
$d_\eta(C[b])=1$ if $C[b]$ is a terminal label. Thus $d_\eta(C)$ is the maximal
number of node/terminal layers on a path beginning at $C$. Set
\[
D_\eta:=\max_{C\in\cT_\eta} d_\eta(C)<\infty.
\]
Whenever $b\in B_C$, the remaining depth below $b$ is at most $d_\eta(C)-1$.
Accordingly, boundary continuation values will always be stored in
$[0,d_\eta(C)-1]^N$, and the corresponding local gains will lie in
$[0,d_\eta(C)]$. In particular a cemetery-routing path contains at most
$D_\eta-1$ strict child moves, hence at most $D_\eta$ cemetery exits.
\end{remark}

\begin{assumption}[First-exit controller realization]
\label{ass:controller-input}
Fix $\eta\in(0,1)$ and a node $C\in\cT_\eta$. For every stationary mixed kernel
$\sigma_C\in\Sigma_C$ and every cap $L\ge 1$, the universal
$\eta$-tree representation supplies a public finite-memory first-exit block
controller $\mathcal R_{C,L}^{\sigma_C}$. During every live stage with current
phase-state $\xi\in\Xi_C$, the compliant action profile is sampled from the
prescribed stationary kernel $\sigma_C(\cdot\mid\xi)$; the additional public
memory is only an implementation device and does not alter the one-step
coefficients on $\Xi_C$. The controller has the following properties.
\begin{enumerate}[label=\textnormal{(R\arabic*)},leftmargin=2.8em]
\item If the block starts from an entry phase-state in $\Xi_C$, its live path
evolves on the genuine augmented state space
$\Xi_C^{\mathrm{aug}}=\Xi_C\sqcup\partial B_C$ until the stopping time
\[
\tau_{C,L}
:=
\inf\{t\ge 1:X_t\in\partial B_C\}\wedge L.
\]
The block consumes exactly $\tau_{C,L}$ stages of the game clock. The parameter
$L$ is only a truncation cap.
\item At most one boundary label is emitted. If $\tau_{C,L}<L$ and
$X_{\tau_{C,L}}=\partial_b$, then that unique label is $b$.
\item Let $N_{C,L}(\xi,a)$ be the number of live stages
$t\in\{0,\dots,\tau_{C,L}-1\}$ at which $(X_t,A_t)=(\xi,a)$, and let
$I_{C,L}(b):=\1_{\{X_{\tau_{C,L}}=\partial_b\}}$. Define the empirical live-time
rates
\[
\widehat\mu_{C,L}^{\mathrm{rate}}(\xi,a):=\frac{N_{C,L}(\xi,a)}{\tau_{C,L}},
\qquad
\widehat\beta_{C,L}(b):=\frac{I_{C,L}(b)}{\tau_{C,L}}.
\]
Then, uniformly over entry states, there exists $K_C(\sigma_C)<\infty$ such that
\begin{equation}
\label{eq:controller-input-occ}
\E\Bigl[
\bigl\|(\widehat\mu_{C,L}^{\mathrm{rate}},\widehat\beta_{C,L})
-
\bar\mu_C^\sigma
\bigr\|_1
\Bigr]
\le \frac{K_C(\sigma_C)}{\sqrt L}.
\end{equation}
\item For every continuation vector
$v_C(\cdot)\in[0,d_\eta(C)-1]^{B_C\times N}$, the controller's total augmented
block reward is
\[
\widetilde U_{i,C,L}
:=
\sum_{t=0}^{\tau_{C,L}-1} u_i^t
+
\sum_{b\in B_C}I_{C,L}(b)\,v_{i,C}(b).
\]
\item Uniformly over entry states and players,
\begin{equation}
\label{eq:controller-input-renewal}
\left|
\frac{\E[\widetilde U_{i,C,L}]}{\E[\tau_{C,L}]}
-
\left(
\sum_{(\xi,a)\in\Omega_C}\mu_C^\sigma(\xi,a)u_i(\xi,a)
+
\sum_{b\in B_C}\beta_C^\sigma(b)v_{i,C}(b)
\right)
\right|
\le \frac{K_C(\sigma_C)}{\sqrt L}.
\end{equation}
\end{enumerate}
\end{assumption}


\subsection{Pure and mixed stationary kernels inside a node}

Fix a node $C$. A \emph{pure stationary selector} of player $i$ in node $C$ is a
map
\[
f_i: \Xi_C\to \bigsqcup_{\xi\in\Xi_C}A_i(\xi)
\quad\text{with }f_i(\xi)\in A_i(\xi).
\]
The finite set of such selectors is denoted by $F_{i,C}$. A \emph{stationary
mixed kernel} of player $i$ is an element of
\[
\Sigma_{i,C}:=\prod_{\xi\in\Xi_C}\Delta(A_i(\xi)).
\]
We write $\Sigma_C:=\prod_{i\in N}\Sigma_{i,C}$.

\begin{remark}
A point of $\Sigma_{i,C}$ may be viewed either as a statewise mixed-action kernel
or, equivalently, as a product distribution over the finite selector set
$F_{i,C}$. We work with statewise mixed kernels because the deviation MDP
coefficients are affine in those kernels; this is the precise form needed later.
The local equilibrium existence result in Section~\ref{sec:local} is equivalent to
Nash's theorem for the finite selector game.
\end{remark}

\section{The bookkeeping node system and the linear exit map}
\label{sec:bookkeeping}

\subsection{Bookkeeping kernel and balance equations}

Fix a node $C\in\cT_\eta$. This section does \emph{not} describe the executed
block dynamics. It records only the internal bookkeeping chain used to define the
convex occupation package. Exit mass is instantaneously re-injected into internal
states through $\rho_C$, so the bookkeeping chain lives on $\Xi_C$ alone.

For a stationary mixed kernel $\sigma_C\in\Sigma_C$, define the induced internal
kernel
\[
Q_C^{\sigma}(\xi'\mid \xi)
:=
\sum_{a\in\prod_j A_j(\xi)} \sigma_C(a\mid \xi)
P_C^0(\xi'\mid \xi,a),
\]
and the induced exit kernel
\[
\Lambda_C^{\sigma}(b\mid\xi)
:=
\sum_{a\in\prod_j A_j(\xi)}\sigma_C(a\mid\xi)\lambda_C(b\mid\xi,a).
\]
The corresponding bookkeeping kernel on $\Xi_C$ is
\begin{equation}
\label{eq:bookkeeping-kernel}
\widetilde Q_C^{\sigma}(\xi'\mid \xi)
=
Q_C^{\sigma}(\xi'\mid \xi)
+
\sum_{b\in B_C}\Lambda_C^{\sigma}(b\mid \xi)\rho_C(\xi'\mid b).
\end{equation}
Its row sums are equal to one by construction.

If $\pi_C^{\sigma,\mathrm{book}}$ denotes an invariant law of
$\widetilde Q_C^{\sigma}$, define the internal occupation measure and exit rate by
\[
\mu_C^{\sigma}(\xi,a):=\pi_C^{\sigma,\mathrm{book}}(\xi)\sigma_C(a\mid\xi),
\qquad
\beta_C^{\sigma}(b):=\sum_{(\xi,a)\in\Omega_C}
\mu_C^{\sigma}(\xi,a)\lambda_C(b\mid\xi,a).
\]
The pair is denoted by
\[
\bar\mu_C^{\sigma}=(\mu_C^{\sigma},\beta_C^{\sigma}).
\]


\begin{definition}
\label{def:PC}
The \emph{internal feasible set} of node $C$ is the set of all
$\mu\in\R_+^{\Omega_C}$ for which there exists $\beta\in\R_+^{B_C}$ such that
\begin{align}
\sum_{(\xi,a)\in\Omega_C}\mu(\xi,a)&=1,
\label{eq:norm-PC}
\\
\pi_\mu(\xi')
&=
\sum_{(\xi,a)\in\Omega_C}\mu(\xi,a)P_C^0(\xi'\mid\xi,a)
+
\sum_{b\in B_C}\beta(b)\rho_C(\xi'\mid b)
\qquad (\xi'\in\Xi_C),
\label{eq:internal-balance}
\\
\beta(b)
&=
\sum_{(\xi,a)\in\Omega_C}\mu(\xi,a)\lambda_C(b\mid\xi,a)
\qquad (b\in B_C),
\label{eq:beta-balance}
\end{align}
where
\[
\pi_\mu(\xi):=\sum_{a\in\prod_j A_j(\xi)}\mu(\xi,a).
\]
We denote this set by $\cP_C$.
\end{definition}


\begin{remark}
The restart kernel $\rho_C$ belongs only to the bookkeeping balance equations
\eqref{eq:internal-balance}. It does not describe what the executed block does
after leaving node $C$. Executed blocks terminate at the first cemetery state and
are routed globally in Section~\ref{sec:realization}.
\end{remark}

\subsection{The graph theorem for exit fluxes}

Define the linear exit map
\begin{equation}
\label{eq:LC-def}
L_C:\R^{\Omega_C}\to \R^{B_C},
\qquad
(L_C\mu)(b):=
\sum_{(\xi,a)\in\Omega_C}\mu(\xi,a)\lambda_C(b\mid\xi,a).
\end{equation}
The associated restart operator is
\begin{equation}
\label{eq:RC-def}
R_C:\R^{B_C}\to \R^{\Xi_C},
\qquad
(R_C\beta)(\xi'):=\sum_{b\in B_C}\beta(b)\rho_C(\xi'\mid b).
\end{equation}


\begin{proposition}
\label{prop:PC-polytope}
For every node $C$, the set $\cP_C$ is a nonempty compact convex polytope.
Moreover, $\mu\in\cP_C$ if and only if
\begin{align}
\sum_{(\xi,a)}\mu(\xi,a)&=1,
\label{eq:PC-short1}
\\
\pi_\mu &= P_C^0\mu + R_C L_C\mu,
\label{eq:PC-short2}
\end{align}
where $(P_C^0\mu)(\xi'):=\sum_{(\xi,a)}\mu(\xi,a)P_C^0(\xi'\mid\xi,a)$.
\end{proposition}

\begin{proof}
The defining relations in \Cref{def:PC} are finitely many linear equalities and
inequalities in the variables $(\mu,\beta)$. Hence the feasible set in
$\R_+^{\Omega_C}\times\R_+^{B_C}$ is a compact convex polyhedron; compactness
comes from \eqref{eq:norm-PC}. Projecting onto the $\mu$-coordinates shows that
$\cP_C$ is a compact convex polytope.

If $\mu\in\cP_C$, the auxiliary vector $\beta$ is unique by
\eqref{eq:beta-balance}, namely $\beta=L_C\mu$. Substituting this into
\eqref{eq:norm-PC} and \eqref{eq:internal-balance} gives
\eqref{eq:PC-short1}--\eqref{eq:PC-short2}. Conversely, if
\eqref{eq:PC-short1}--\eqref{eq:PC-short2} hold and one sets $\beta=L_C\mu$, then
all three relations from \Cref{def:PC} are satisfied.
\end{proof}


\begin{definition}
\label{def:MCbook}
The \emph{bookkeeping augmented feasible set} of node $C$ is
\[
\cM_C^{\mathrm{book}}
:=
\{(\mu,\beta)\in\cP_C\times\R_+^{B_C}:\beta=L_C\mu\}.
\]
The affine graph map is
\[
J_C:\cP_C\to \cM_C^{\mathrm{book}},
\qquad J_C(\mu)=(\mu,L_C\mu).
\]
\end{definition}

\begin{theorem}[Graph theorem for the bookkeeping node]
\label{thm:graph-theorem}
For every node $C$,
\[
\cM_C^{\mathrm{book}} = J_C(\cP_C).
\]
The map $J_C$ is affine and injective. In particular the exit coordinates are not
independent variables: once the internal part $\mu$ is fixed, the exit mass is
uniquely determined as $L_C\mu$.
\end{theorem}

\begin{proof}
The identity is immediate from \Cref{def:MCbook}. Affineness follows from the
linearity of $L_C$, and injectivity is obvious from the first coordinate.
\end{proof}

\begin{corollary}
\label{cor:no-vp3}
There is no separate compatibility condition between internal occupation and exit
mass. The feasible bookkeeping set is exactly the graph of the linear exit map.
\end{corollary}

\begin{proof}
This is a restatement of \Cref{thm:graph-theorem}.
\end{proof}

\subsection{Actual bookkeeping packages induced by stationary mixed kernels}

\begin{definition}
\label{def:actual-local-package}
For a stationary mixed kernel $\sigma_C\in\Sigma_C$, its \emph{actual bookkeeping
occupation} is the pair $\bar\mu_C^{\sigma}$ induced by an invariant law of the
bookkeeping kernel \eqref{eq:bookkeeping-kernel}. By construction,
$\bar\mu_C^{\sigma}\in \cM_C^{\mathrm{book}}$.
\end{definition}

\begin{proof}
Let $\pi=\pi_C^{\sigma,\mathrm{book}}$ be an invariant law of
$\widetilde Q_C^{\sigma}$. Writing the stationarity equation for
$\widetilde Q_C^{\sigma}$ and expanding \eqref{eq:bookkeeping-kernel} gives
exactly \eqref{eq:norm-PC}--\eqref{eq:beta-balance} with
$\mu=\mu_C^{\sigma}$ and $\beta=\beta_C^{\sigma}$.
\end{proof}

\section{Local average-reward games and Bellman certificates}
\label{sec:local}

\subsection{Local continuation values}

Fix $\eta\in(0,1)$ and a node $C$. Let $x$ be a package of continuation data on the
whole tree. For each boundary state $b\in B_C$ the value
$v_C^x(b)\in[0,d_\eta(C)-1]^N$ is read from the continuation coordinates of $x$:
if $C[b]$ is a child node $D$, then $v_C^x(b)$ is the continuation vector stored
at $b$ in $x$ and required later to be $\eta$-close to $g_D$; if $C[b]$ is a
terminal label, then $v_C^x(b)=w(b)$ is the payoff vector of the terminal
continuation profile attached to that label. Throughout this section $x$ is
considered fixed.

\subsection{The genuine augmented kernel and the augmented stage payoff}

Fix a node $C$ and a stationary mixed kernel $\sigma_C\in\Sigma_C$. The
\emph{genuine augmented state space} is
\[
\Xi_C^{\mathrm{aug}}:=\Xi_C\sqcup \partial B_C,
\qquad
\partial B_C:=\{\partial_b:b\in B_C\}.
\]
The executed augmented kernel is
\begin{equation}
\label{eq:aug-kernel}
P_C^{\sigma,\mathrm{aug}}(z\mid y)
=
\begin{cases}
Q_C^{\sigma}(\xi'\mid \xi), & y=\xi\in\Xi_C,\ z=\xi'\in\Xi_C,\\[0.3em]
\Lambda_C^{\sigma}(b\mid \xi), & y=\xi\in\Xi_C,\ z=\partial_b,\\[0.3em]
1, & y=z=\partial_b,\\[0.3em]
0, & \text{otherwise.}
\end{cases}
\end{equation}
Thus cemetery states are absorbing. The executed block controller will terminate at
the first cemetery hit; the absorbing completion in \eqref{eq:aug-kernel} is used
only to make the kernel stochastic.

For a fixed player $i$ and opponents' kernel $\sigma_{-i}$, the stage reward part,
the internal transition part, and the exit part are
\begin{align}
\label{eq:r-sigma}
r_{i,C}^{\sigma_{-i}}(\xi,a_i)
&:=
\sum_{a_{-i}}\sigma_{-i}(a_{-i}\mid\xi)
\,u_i\bigl(\xi,(a_i,a_{-i})\bigr),
\\
\label{eq:Q-sigma}
Q_{i,C}^{\sigma_{-i}}(\xi'\mid \xi,a_i)
&:=
\sum_{a_{-i}}\sigma_{-i}(a_{-i}\mid\xi)
\,P_C^0\bigl(\xi'\mid\xi,(a_i,a_{-i})\bigr),
\\
\label{eq:q-exit-sigma}
\Lambda_{i,C}^{\sigma_{-i}}(b\mid\xi,a_i)
&:=
\sum_{a_{-i}}\sigma_{-i}(a_{-i}\mid\xi)
\,\lambda_C\bigl(b\mid\xi,(a_i,a_{-i})\bigr).
\end{align}
The associated \emph{augmented one-step reward} is
\begin{equation}
\label{eq:augmented-reward}
\widetilde r_{i,C}^{\sigma_{-i},x}(\xi,a_i)
:=
r_{i,C}^{\sigma_{-i}}(\xi,a_i)
+
\sum_{b\in B_C}\Lambda_{i,C}^{\sigma_{-i}}(b\mid\xi,a_i)v_{i,C}^x(b).
\end{equation}
This is exactly the expected payoff of one block step when the block is still in
$\Xi_C$: the ordinary stage payoff is earned inside $\Xi_C$, and if the next state
is the first cemetery state $\partial_b$, the controller records the continuation
value $v_{i,C}^x(b)$ at that step.

The Bellman operator is therefore
\begin{equation}
\label{eq:Bellman-operator}
(T_{i,C}^{\sigma_{-i},x}h)(\xi)
:=
\max_{a_i\in A_i(\xi)}\Bigl[
 r_{i,C}^{\sigma_{-i}}(\xi,a_i)
 +\sum_{\xi'\in\Xi_C}Q_{i,C}^{\sigma_{-i}}(\xi'\mid\xi,a_i)h(\xi')
 +\sum_{b\in B_C}\Lambda_{i,C}^{\sigma_{-i}}(b\mid\xi,a_i)v_{i,C}^x(b)
 \Bigr].
\end{equation}
Equivalently, this is the Bellman operator of the augmented MDP on
$\Xi_C^{\mathrm{aug}}$ with cemetery biases fixed at zero.

\begin{lemma}
\label{lem:coeff-affine}
For fixed $C$, $i$, and $x$, the maps
\[
\sigma_{-i}\mapsto r_{i,C}^{\sigma_{-i}},
\qquad
\sigma_{-i}\mapsto Q_{i,C}^{\sigma_{-i}},
\qquad
\sigma_{-i}\mapsto \Lambda_{i,C}^{\sigma_{-i}},
\qquad
\sigma_{-i}\mapsto \widetilde r_{i,C}^{\sigma_{-i},x}
\]
are affine. Consequently there exists a finite constant $K_{i,C}$ such that for
all $\sigma_{-i},\tau_{-i}$,
\begin{align}
\|\widetilde r_{i,C}^{\tau_{-i},x}-\widetilde r_{i,C}^{\sigma_{-i},x}\|_\infty
&\le K_{i,C}\|\tau_{-i}-\sigma_{-i}\|_1,
\label{eq:r-lip}
\\
\sup_{\xi,a_i}\|Q_{i,C}^{\tau_{-i}}(\cdot\mid\xi,a_i)-Q_{i,C}^{\sigma_{-i}}(\cdot\mid\xi,a_i)\|_{\mathrm{TV}}
&\le K_{i,C}\|\tau_{-i}-\sigma_{-i}\|_1,
\label{eq:Q-lip}
\\
\sup_{\xi,a_i}\|\Lambda_{i,C}^{\tau_{-i}}(\cdot\mid\xi,a_i)-\Lambda_{i,C}^{\sigma_{-i}}(\cdot\mid\xi,a_i)\|_{1}
&\le K_{i,C}\|\tau_{-i}-\sigma_{-i}\|_1.
\label{eq:Lambda-lip}
\end{align}
\end{lemma}

\begin{proof}
Each coefficient is a finite linear combination of the components of
$\sigma_{-i}(\cdot\mid\xi)$, with coefficients given by the primitive reward,
transition, and continuation data. This proves affineness. The Lipschitz bounds
follow because all underlying arrays are bounded and finite.
\end{proof}

\subsection{Local best responses}

We now formulate the local best-response correspondence. Since the local state and
action sets are finite, every unilateral local problem is a finite average-reward
MDP with communicating internal part and bounded terminal continuation rewards.

\begin{definition}
\label{def:local-BR-set}
For fixed $x$, node $C$, and player $i$, let $\mathfrak B_{i,C}(x)$ be the set of
triples $(\sigma_{i,C},g_{i,C},h_{i,C})$ such that
\begin{enumerate}[label=\textnormal{(BR\arabic*)},leftmargin=2.9em]
\item $\sigma_{i,C}\in \Sigma_{i,C}$;
\item $g_{i,C}\in[0,d_\eta(C)]$ and
$h_{i,C}\in[-B_\eta^0,B_\eta^0]^{\Xi_C}$;
\item $h_{i,C}$ is normalized by
\[
h_{i,C}(\xi_C^\circ)=0;
\]
\item for every $\xi\in\Xi_C$,
\begin{equation}
\label{eq:Bellman-ineq-def}
g_{i,C}+h_{i,C}(\xi)
\ge
r_{i,C}^{\sigma_{-i,C}^x}(\xi,a_i)
+
\sum_{\xi'}Q_{i,C}^{\sigma_{-i,C}^x}(\xi'\mid\xi,a_i)h_{i,C}(\xi')
+
\sum_b \Lambda_{i,C}^{\sigma_{-i,C}^x}(b\mid\xi,a_i)v_{i,C}^x(b)
\end{equation}
for every $a_i\in A_i(\xi)$;
\item equality holds in \eqref{eq:Bellman-ineq-def} for every action $a_i$ in the
support of $\sigma_{i,C}(\cdot\mid\xi)$.
\end{enumerate}
\end{definition}

\begin{proposition}
\label{prop:local-BR}
Assume condition~\textnormal{(a)}. For every fixed $x$, node $C$, and player $i$,
the set $\mathfrak B_{i,C}(x)$ is nonempty, compact, and convex. The
correspondence $x\mapsto \mathfrak B_{i,C}(x)$ has closed graph.
\end{proposition}

\begin{proof}
Fix $x$, $C$, and $i$. The Bellman system
\eqref{eq:Bellman-ineq-def} is exactly the dual optimality system for the finite
average-reward MDP faced by player $i$ against the frozen environment encoded by
$x$; see \Cref{prop:mdp-duality}. Because the frozen continuation vector lies in
$[0,d_\eta(C)-1]^N$ and the stage payoffs lie in $[0,1]$, the augmented one-step
reward in that MDP is bounded by $d_\eta(C)$. Hence there exists an optimal
stationary mixed kernel together with a gain-bias pair satisfying
\eqref{eq:Bellman-ineq-def} and the support equalities, with
$g_{i,C}\in[0,d_\eta(C)]$. Therefore $\mathfrak B_{i,C}(x)$ is nonempty.

Compactness follows because $\Sigma_{i,C}$ is a product of simplices and the gain
and bias coordinates are confined to compact boxes. The defining equalities and
inequalities are closed.

For convexity, condition~\textnormal{(a)} implies that the internal part of the
deviation MDP is communicating on $\Xi_C$. For communicating finite MDPs, the
optimal bias is unique up to an additive constant; after the linear normalization
$h(\xi_C^\circ)=0$, the optimal bias is unique, see
\Cref{prop:bias-unique} and \cite[Chapter~8]{Puterman1994}. The same is true for
the optimal gain. Hence, once $x$ is fixed, there is a unique normalized optimal
pair $(g^*,h^*)$. A stationary mixed kernel is optimal if and only if, at every
state, it assigns mass only to actions attaining equality in the Bellman system
for $(g^*,h^*)$. The set of such kernels is a product of simplices and hence
convex. Thus $\mathfrak B_{i,C}(x)$ is convex.

For the graph property, let $x^m\to x$ and let
$(\sigma_i^m,g_i^m,h_i^m)\in \mathfrak B_{i,C}(x^m)$ converge to
$(\sigma_i,g_i,h_i)$. The coefficient arrays from \Cref{lem:coeff-affine} depend
continuously on $x$, hence passing to the limit in the Bellman equalities and
inequalities shows $(\sigma_i,g_i,h_i)\in\mathfrak B_{i,C}(x)$.
\end{proof}

\subsection{Local equilibrium existence}

\begin{proposition}[Local stationary Nash equilibrium]
\label{prop:local-equilibrium}
Assume condition~\textnormal{(a)}. Fix $x$ and a node $C$. There exists a
stationary mixed kernel $\sigma_C^*\in\Sigma_C$ such that for every player $i$,
$(\sigma_{i,C}^*,g_{i,C}^*,h_{i,C}^*)\in \mathfrak B_{i,C}(x)$ for some
$(g_{i,C}^*,h_{i,C}^*)$. Equivalently, the local average game at node $C$ with
continuation values frozen at $x$ has a stationary Nash equilibrium, and every
player admits a Bellman certificate at that equilibrium.
\end{proposition}

\begin{proof}
Fix the continuation vector $v_C^x(\cdot)$ extracted from $x$. For each player
$i$ and every opponents' kernel
$\sigma_{-i,C}\in\prod_{j\ne i}\Sigma_{j,C}$, let
\[
\BR_{i,C}^x(\sigma_{-i,C})
:=
\{\tau_{i,C}\in\Sigma_{i,C}: \exists g,h\text{ with }
(\tau_{i,C},g,h)\text{ satisfying \eqref{eq:Bellman-ineq-def} against }
\sigma_{-i,C}\text{ and }v_C^x\}.
\]
By the same MDP-duality argument as in \Cref{prop:local-BR}, each value
$\BR_{i,C}^x(\sigma_{-i,C})$ is nonempty, compact, and convex, and the graph of
the playerwise best-response correspondence is closed. Define the product
correspondence on $\Sigma_C$ by
\[
\BR_C^x(\sigma_C):=\prod_{i\in N}\BR_{i,C}^x(\sigma_{-i,C}).
\]
Its values are therefore nonempty, compact, and convex, and its graph is closed.
Since $\Sigma_C$ is a compact convex product of simplices, Kakutani's theorem
\cite{Kakutani1941} yields a fixed point $\sigma_C^*\in\BR_C^x(\sigma_C^*)$. For
each player $i$, choose a Bellman certificate $(g_{i,C}^*,h_{i,C}^*)$ associated
with the best-response relation at that fixed point. This is exactly the required
local stationary Nash package.
\end{proof}

\begin{remark}
This proposition is the correct multiplayer replacement for the v4 appendix, which
contained only one-player LP duality. One-player LP duality is still needed, but
only to identify each player's best-response set. Existence of a simultaneous
local Nash package requires the fixed-point step above.
\end{remark}

\section{The global correspondence with frozen continuation values}
\label{sec:global}

\subsection{Ambient package space}

A global package contains, at each node and for each player, a stationary mixed
kernel and a Bellman certificate, together with nodewise continuation values.
There is \emph{no} self-consistency requirement in the ambient space.

\begin{definition}
\label{def:Aeta}
For fixed $\eta\in(0,1)$, the ambient package space $\cA_\eta$ consists of all
families
\[
x=\Bigl((\sigma_{i,C},g_{i,C},h_{i,C})_{i\in N},\ (v_C(b))_{b\in B_C}\Bigr)_{C\in\cT_\eta}
\]
such that:
\begin{enumerate}[label=\textnormal{(A\arabic*)},leftmargin=2.8em]
\item $\sigma_{i,C}\in \Sigma_{i,C}$ for all $i,C$;
\item $g_{i,C}\in[0,d_\eta(C)]$ for all $i,C$;
\item $h_{i,C}\in[-B_\eta^0,B_\eta^0]^{\Xi_C}$ and
$h_{i,C}(\xi_C^\circ)=0$ for all $i,C$;
\item $v_C(b)\in[0,d_\eta(C)-1]^N$ for all $C$ and $b\in B_C$.
\end{enumerate}
\end{definition}

\begin{proposition}
\label{prop:Aeta-compact}
For every $\eta\in(0,1)$, the set $\cA_\eta$ is compact and convex.
\end{proposition}

\begin{proof}
Each $\Sigma_{i,C}$ is a product of simplices, hence compact and convex. The gain
coordinates lie in the depth-dependent intervals $[0,d_\eta(C)]$, the bias
coordinates lie in a compact box intersected with the linear normalization
hyperplane $h(\xi_C^\circ)=0$, and the continuation values lie in the
depth-dependent cubes $[0,d_\eta(C)-1]^N$. Because the tree is finite,
$\cA_\eta$ is a finite product of compact convex sets.
\end{proof}

\subsection{Global best response}

The key point is that continuation values are frozen from the reference package
$x$. The response package $y$ is allowed to differ from $x$.

\begin{definition}
\label{def:global-BR}
For $x\in\cA_\eta$, define $\BR_\eta(x)$ to be the set of all packages
$y\in\cA_\eta$ such that:
\begin{enumerate}[label=\textnormal{(G\arabic*)},leftmargin=2.8em]
\item for every node $C$ and every player $i$,
\[
(\sigma_{i,C}^y,g_{i,C}^y,h_{i,C}^y)\in \mathfrak B_{i,C}(x);
\]
\item whenever $b\in B_C$ leads to a child node $D=C[b]$, one has
\begin{equation}
\label{eq:frozen-continuation}
\|v_C^y(b)-g_D^x\|_\infty\le \eta,
\end{equation}
where $g_D^x=(g_{i,D}^x)_{i\in N}$;
\item if $b\in B_C$ leads to a terminal label with continuation vector
$w(b)\in[0,1]^N\subseteq[0,d_\eta(C)-1]^N$, then
\[
v_C^y(b)=w(b).
\]
\end{enumerate}
\end{definition}

\begin{remark}
The response package $y$ is assembled against \emph{frozen} continuation values and
frozen opponents from $x$. There is no circularity here. One does not ask $y$ to
be self-consistent while defining $\BR_\eta(x)$; that happens only at a fixed
point. This is exactly the continuation repair already made in v5 and retained in
v6.
\end{remark}

\begin{theorem}[Global Kakutani fixed point]
\label{thm:kakutani-global}
Assume condition~\textnormal{(a)}. For every $\eta\in(0,1)$, the correspondence
$\BR_\eta:\cA_\eta\rightrightarrows\cA_\eta$ has nonempty compact convex values and
closed graph. Consequently it has a fixed point $x^*\in\cA_\eta$.
\end{theorem}

\begin{proof}
By \Cref{prop:Aeta-compact}, the domain $\cA_\eta$ is compact and convex.

Fix $x\in\cA_\eta$. For every node $C$ and player $i$, the set
$\mathfrak B_{i,C}(x)$ is nonempty, compact, and convex by
\Cref{prop:local-BR}. The continuation constraints
\eqref{eq:frozen-continuation} are closed convex constraints, and terminal
continuation coordinates are fixed singleton values. Therefore $\BR_\eta(x)$ is a
finite product of nonempty compact convex sets, hence itself nonempty, compact, and
convex.

For the graph property, let $x^m\to x$ and $y^m\to y$ with
$y^m\in\BR_\eta(x^m)$ for all $m$. For every node $C$ and player $i$ we have
$(\sigma_{i,C}^{y^m},g_{i,C}^{y^m},h_{i,C}^{y^m})\in\mathfrak B_{i,C}(x^m)$. By
the closed-graph property from \Cref{prop:local-BR}, the limit belongs to
$\mathfrak B_{i,C}(x)$. Passing to the limit in the continuation inequalities gives
\eqref{eq:frozen-continuation} for $y$ against $x$. Hence $y\in\BR_\eta(x)$, so
$\BR_\eta$ has closed graph.

Closed graph plus compact values on a compact domain imply upper
hemicontinuity. Kakutani's theorem \cite{Kakutani1941} yields a fixed point.
\end{proof}

\begin{corollary}
\label{cor:local-nash-at-fixed-point}
If $x^*\in\BR_\eta(x^*)$, then for every node $C$ the stationary mixed kernel
$\sigma_C^*:=(\sigma_{i,C}^*)_{i\in N}$ is a local stationary Nash equilibrium of
the node game with continuation values $v_C^{x^*}(\cdot)$, and each player has a
Bellman certificate $(g_{i,C}^*,h_{i,C}^*)$.
\end{corollary}

\begin{proof}
At a fixed point,
$(\sigma_{i,C}^*,g_{i,C}^*,h_{i,C}^*)\in\mathfrak B_{i,C}(x^*)$ for every $i,C$.
Thus each player's stationary mixed kernel is a best response to the others inside
node $C$, against the continuation values carried by $x^*$. That is exactly the
local stationary Nash property.
\end{proof}

\section{Local controllers and global realization}
\label{sec:realization}


\subsection{The node controller and first-exit blocks}

Fix a fixed point $x^*\in\BR_\eta(x^*)$. For each node $C$, write
\[
\sigma_C^*:=(\sigma_{i,C}^*)_{i\in N},
\qquad
\bar\mu_C^*:=(\mu_C^{\sigma_C^*},\beta_C^{\sigma_C^*}).
\]
The local controller for node $C$ is the first-exit controller
$\mathcal R_{C,L}^{\sigma_C^*}$ supplied by \Cref{ass:controller-input}. When the
controller is launched with cap $L$, it evolves on
$\Xi_C^{\mathrm{aug}}=\Xi_C\sqcup\partial B_C$ until the actual stopping time
$\tau_{C,L}\le L$. The cap $L$ controls realization accuracy, but the block
consumes only the \emph{actual} time $\tau_{C,L}$.

For a realized block $B$ in node $C$, let $L(B)$ denote its cap and let
$|B|:=\tau(B)$ denote its actual live duration. Let
$N_B(\xi,a)$ be the number of live stages in $B$ at which $(X_t,A_t)=(\xi,a)$, and
let $I_B(b)$ be the indicator that the block exits at $\partial_b$. The empirical
live-time rates are
\[
\widehat\mu_B^{\mathrm{rate}}(\xi,a):=\frac{N_B(\xi,a)}{|B|},
\qquad
\widehat\beta_B(b):=\frac{I_B(b)}{|B|}.
\]
The total augmented block reward is
\begin{equation}
\label{eq:aug-block-payoff}
\widetilde U_i(B)
:=
\sum_{t\in B^{\mathrm{live}}} u_i^t
+
\sum_{b\in B_C}I_B(b)\,v_{i,C}^{x^*}(b).
\end{equation}
If the block ends at its cap without a cemetery hit, the second term is zero.

\begin{theorem}[Local controller lemma]
\label{thm:local-controller}
Assume \Cref{ass:controller-input}. Fix a node $C$ and the fixed-point kernel
$\sigma_C^*$. Then there exists a constant $K_C<\infty$ such that for every cap
$L\ge 1$, every entry phase-state, and every player $i$,
\begin{equation}
\label{eq:local-controller-occ}
\E\bigl[\|(\widehat\mu_{C,L}^{\mathrm{rate}},\widehat\beta_{C,L})-\bar\mu_C^*\|_1\bigr]
\le \frac{K_C}{\sqrt L},
\end{equation}
and
\begin{equation}
\label{eq:local-controller-payoff}
\left|
\frac{\E[\widetilde U_{i,C,L}]}{\E[\tau_{C,L}]}
-
g_{i,C}^*
\right|
\le \frac{K_C}{\sqrt L}.
\end{equation}
Consequently, for every $\delta>0$ there exists $L_C(\delta)$ such that the block
controller is $\delta$-accurate in the sense of
\eqref{eq:local-controller-occ}--\eqref{eq:local-controller-payoff} whenever the
cap $L$ satisfies $L\ge L_C(\delta)$.
\end{theorem}

\begin{proof}
The occupation estimate \eqref{eq:local-controller-occ} is exactly the first-exit
realization estimate from \Cref{ass:controller-input}, with
$K_C:=K_C(\sigma_C^*)$.

Define the bounded linear reward-rate functional
\[
\Phi_{i,C}(\mu,\beta)
:=
\sum_{(\xi,a)\in\Omega_C}\mu(\xi,a)u_i(\xi,a)
+
\sum_{b\in B_C}\beta(b)v_{i,C}^{x^*}(b).
\]
By \eqref{eq:controller-input-renewal},
\[
\left|
\frac{\E[\widetilde U_{i,C,L}]}{\E[\tau_{C,L}]}
-
\Phi_{i,C}(\bar\mu_C^*)
\right|
\le \frac{K_C'}{\sqrt L}.
\]

It remains to identify the target. Integrating the Bellman equalities for actions
in the support of $\sigma_{i,C}^*(\cdot\mid\xi)$ against the bookkeeping rate
package $\bar\mu_C^*=(\mu_C^*,\beta_C^*)$ and using
\eqref{eq:internal-balance}--\eqref{eq:beta-balance} yields
\[
g_{i,C}^*
=
\sum_{(\xi,a)}\mu_C^*(\xi,a)u_i(\xi,a)
+
\sum_{b\in B_C}\beta_C^*(b)v_{i,C}^{x^*}(b)
-
\sum_{\xi'\in\Xi_C}(R_C\beta_C^*)(\xi')h_{i,C}^*(\xi').
\]
Because the bookkeeping restart is fixed at the reference state,
$\rho_C(\cdot\mid b)=\delta_{\xi_C^\circ}$ for every $b\in B_C$, one has
\[
(R_C\beta_C^*)(\xi')
=
\Bigl(\sum_{b\in B_C}\beta_C^*(b)\Bigr)\delta_{\xi_C^\circ}(\xi').
\]
Hence the restart-bias term vanishes:
\[
\sum_{\xi'\in\Xi_C}(R_C\beta_C^*)(\xi')h_{i,C}^*(\xi')
=
\Bigl(\sum_{b\in B_C}\beta_C^*(b)\Bigr)h_{i,C}^*(\xi_C^\circ)
=
0,
\]
because $h_{i,C}^*(\xi_C^\circ)=0$ by normalization. Therefore
\[
g_{i,C}^*
=
\sum_{(\xi,a)}\mu_C^*(\xi,a)u_i(\xi,a)
+
\sum_{b\in B_C}\beta_C^*(b)v_{i,C}^{x^*}(b)
=
\Phi_{i,C}(\bar\mu_C^*).
\]
Combining the two displays proves \eqref{eq:local-controller-payoff} after
absorbing constants into $K_C$. This is precisely the renewal-reward identity: the
correct normalization uses the actual first-exit time $\tau_{C,L}$, not the cap
$L$.
\end{proof}

\subsection{Frozen superexponential schedule}


\begin{definition}
A sequence $(M_r)_{r\ge 1}$ of positive integers is a
\emph{superexponential schedule} if $M_{r+1}\ge M_r^2$ for all $r$. Its
\emph{freezing at level $R$} is the sequence $(M_r^{(R)})_{r\ge 1}$ defined by
\[
M_r^{(R)}=
\begin{cases}
M_r, & r<R,\\
M_R, & r\ge R.
\end{cases}
\]
We write $M_*:=M_R$ for the tail block length.
\end{definition}

\begin{lemma}
\label{lem:frozen-memory}
A frozen superexponential schedule uses only finite public memory.
\end{lemma}

\begin{proof}
Only the finitely many lengths $M_1,\dots,M_R$ appear. To implement the schedule
one stores the current node, the current position inside the current block, the
current local memory state of the node controller, and the block index truncated at
$R$. All of these coordinates range over finite sets.
\end{proof}


\subsection{Public routing and realization}

Recall from \Cref{rem:tree-depth} that $D_\eta<\infty$ is the maximal number of
node/terminal layers on a root-to-leaf path of the finite universal tree
$\cT_\eta$.

\begin{proposition}[Public routing]
\label{prop:public-routing}
Fix a fixed point $x^*$ and a frozen superexponential schedule. There exists a
public finite-memory profile $\sigma_{\eta,\delta,M_*}$ with the following
properties.
\begin{enumerate}[label=\textnormal{(R\arabic*)},leftmargin=2.8em]
\item At each launched block played in node $C$ with cap $L(B)$, the profile runs
the first-exit controller $\mathcal R_{C,L(B)}^{\sigma_C^*}$ from
\Cref{thm:local-controller}.
\item The block consumes exactly $|B|=\tau(B)\le L(B)$ actual stages and
terminates at its first visit to a cemetery state $\partial_b$, if such a visit
occurs, and otherwise at its cap $L(B)$.
\item If a block terminates at $\partial_b$ and $C[b]$ is a child node, the next
block starts immediately in that child node. If a block terminates at
$\partial_b$ and $C[b]$ is a terminal label, the block schedule stops and play
continues according to the terminal continuation profile attached to $b$, whose
payoff vector is $w(b)$. Since terminal labels have no further tree structure,
no additional blocks are scheduled at that label. If no cemetery state is
visited, the next block starts immediately in the same node.
\item Because routing on cemetery hits always moves to a strict child, every play
path contains at most $D_\eta$ cemetery exits.
\item The public memory stores the current node, the current cap index, the
current position inside the live block, and the current local controller state.
The routing update is publicly observable and deterministic conditional on that
information and on the unique boundary label, if one is emitted.
\end{enumerate}
\end{proposition}

\begin{proof}
Items~(R1) and~(R2) are part of the controller realization in
\Cref{ass:controller-input}. Because the first cemetery hit is unique, at most one
boundary label is emitted. The successor map $b\mapsto C[b]$ belongs to the node
data, so once a label $b$ is observed the next routing decision is deterministic:
if $C[b]$ is a child node, the next block starts there; if $C[b]$ is a terminal
label, control passes to the attached terminal continuation profile and no
further block is scheduled. If no label is emitted, the controller remains in the
same node. The bound by $D_\eta$ follows from finiteness of the tree and the fact
that every cemetery-routing move to a child strictly decreases the remaining
depth. The memory requirements are finite by \Cref{lem:frozen-memory}, and all
relevant quantities are publicly observable.
\end{proof}

\begin{theorem}[Realization theorem]
\label{thm:realization}
Assume condition~\textnormal{(a)} and \Cref{ass:controller-input}. Fix
$\eta\in(0,1)$ and a fixed point $x^*\in\BR_\eta(x^*)$. Let
$\sigma_{\eta,\delta,M_*}$ be the public profile obtained from
\Cref{prop:public-routing} by freezing a superexponential schedule at level $R$.
Then there exist constants $C_\eta,K_\eta<\infty$ such that:
\begin{enumerate}[label=\textnormal{(Q\arabic*)},leftmargin=2.8em]
\item if a complete block $B$ is played in node $C$ and its cap satisfies
$L(B)\ge M_*$, then
\begin{equation}
\label{eq:block-occupation-realization}
\E\bigl[\|(\widehat\mu_B^{\mathrm{rate}},\widehat\beta_B)-\bar\mu_C^*\|_1
\mid \cF_{\mathrm{in}(B)}\bigr]
\le \delta;
\end{equation}
\item for every player $i$ and every complete block $B$ played in node $C$
whose cap satisfies $L(B)\ge M_*$,
\begin{equation}
\label{eq:block-payoff-realization}
\left|
\E\bigl[\widetilde U_i(B)\mid \cF_{\mathrm{in}(B)}\bigr]
-
g_{i,C}^*\,\E\bigl[|B|\mid \cF_{\mathrm{in}(B)}\bigr]
\right|
\le C_\eta\delta\,\E\bigl[|B|\mid \cF_{\mathrm{in}(B)}\bigr];
\end{equation}
\item for every horizon $T$,
\begin{equation}
\label{eq:realization-tail-bound}
\left|
\gamma_i^T(s,\sigma_{\eta,\delta,M_*})
-
\frac1T\sum_{B\subseteq [1,T]} |B|\,g_{i,C(B)}^*
\right|
\le C_\eta\delta + \frac{K_\eta}{T},
\end{equation}
where the sum is over complete blocks contained in $[1,T]$ and $C(B)$ denotes the
node of block $B$.
\end{enumerate}
\end{theorem}

\begin{proof}
Item~(Q1) is \Cref{thm:local-controller} at the chosen tail scale, conditioned on
the entry state of the block.

Item~(Q2) is the renewal-reward estimate from \Cref{thm:local-controller}
applied to tail blocks, again conditioned on the entry state and with constants
enlarged uniformly over the finite tree.

For (Q3), decompose the horizon into the complete blocks contained in $[1,T]$ plus
at most one initial incomplete fragment and one terminal incomplete fragment.
Among the complete blocks, only finitely many belong to the frozen prelude and
have cap $<M_*$. Their combined contribution to the realized payoff and to the
weighted-gain sum is bounded by a finite constant depending only on the frozen
prelude, the tail cap $M_*$, and the finite tree object; absorb that constant
into $K_\eta$. Summing
\eqref{eq:block-payoff-realization} over the remaining tail blocks and using
\[
\sum_{B\subseteq [1,T]} |B|\le T
\]
gives the weighted gain approximation.

It remains to compare augmented and ordinary payoffs. A complete block differs
from its ordinary stage-payoff sum only when it exits to a cemetery state, in
which case the augmented payoff contains one continuation term lying in
$[0,D_\eta]$. By \Cref{prop:public-routing}, the play can contain at most $D_\eta$
such exits. Hence the total augmented/ordinary discrepancy over all complete
blocks inside $[1,T]$ is at most $D_\eta^2$.

The incomplete fragments contribute at most a finite constant depending only on the
frozen prelude, the tail cap $M_*$, and the finite tree object. Absorbing these
constants into $K_\eta$ yields \eqref{eq:realization-tail-bound}.
\end{proof}

\section{Condition \texorpdfstring{$(a)$}{(a)}, bias bounds, and the proof of the main theorem}
\label{sec:proof-main}

\subsection{Condition \texorpdfstring{$(a)$}{(a)} in primitive data}

We now formulate the nodewise irreducibility hypothesis directly from the genuine
augmented transition data.

\begin{condition}[Robust nodewise irreducibility]
\label{cond:a}
For every $\eta\in(0,1)$, every node $C\in\cT_\eta$, every player $i\in N$, every
statewise mixed kernel of the other players
$\sigma_{-i}\in\prod_{j\ne i}\Sigma_{j,C}$, and every pure action selection
$a_i(\cdot)$ of player $i$, the augmented kernel
\[
P_{i,C}^{\sigma_{-i},a_i,\mathrm{aug}}(z\mid y)
\]
on $\Xi_C^{\mathrm{aug}}=\Xi_C\sqcup\partial B_C$, defined by
\[
P_{i,C}^{\sigma_{-i},a_i,\mathrm{aug}}(z\mid y)
=
\begin{cases}
\displaystyle
\sum_{a_{-i}}\sigma_{-i}(a_{-i}\mid\xi)
P_C^0\bigl(\xi'\mid\xi,(a_i(\xi),a_{-i})\bigr),
& y=\xi\in\Xi_C,\ z=\xi'\in\Xi_C,\\[0.7em]
\displaystyle
\sum_{a_{-i}}\sigma_{-i}(a_{-i}\mid\xi)
\lambda_C\bigl(b\mid\xi,(a_i(\xi),a_{-i})\bigr),
& y=\xi\in\Xi_C,\ z=\partial_b,\\[0.7em]
1, & y=z=\partial_b,\\
0, & \text{otherwise,}
\end{cases}
\]
has the property that its restriction to the internal states remains irreducible
before absorption: for every $\xi,\zeta\in\Xi_C$ there exists
$1\le n\le |\Xi_C|-1$ with
\[
\bigl(P_{i,C}^{\sigma_{-i},a_i,\mathrm{aug}}\bigr)^n(\zeta\mid\xi)>0.
\]
\end{condition}

\begin{remark}
Condition~\textnormal{(a)} has a direct operational meaning: no single player can,
by choosing actions state by state while the others keep their local kernels,
create a noncommunicating internal class inside a metastable node before the block
is absorbed at a cemetery state.
\end{remark}

\begin{example}
\label{ex:cond-a-trivial}
Suppose that for every node $C$, every phase-state $\xi\in\Xi_C$, and every full
action profile $a\in\prod_j A_j(\xi)$, the internal transition kernel
$P_C^0(\cdot\mid\xi,a)$ belongs to an irreducible family whose directed support
graph is the same strongly connected graph on $\Xi_C$. Then every convex mixture
of those kernels still gives positive probability to each support edge, and hence
every augmented kernel from \Cref{cond:a} can reach any internal state from any
other before absorption. Therefore condition~\textnormal{(a)} holds.
\end{example}

\subsection{A uniform irreducibility modulus}

For fixed $\eta$, define
\begin{equation}
\label{eq:kappa-eta}
\kappa_\eta
:=
\min_{C\in\cT_\eta}
\min_{i\in N}
\inf_{\sigma_{-i}\in\prod_{j\ne i}\Sigma_{j,C}}
\min_{a_i(\cdot)}
\min_{\xi,\zeta\in\Xi_C}
\max_{1\le n\le |\Xi_C|-1}
\Bigl(P_{i,C}^{\sigma_{-i},a_i,\mathrm{aug}}\Bigr)^n(\zeta\mid\xi).
\end{equation}

\begin{lemma}
\label{lem:kappa-positive}
Assume condition~\textnormal{(a)}. Then $\kappa_\eta>0$ for every $\eta\in(0,1)$.
\end{lemma}

\begin{proof}
Fix $C$, $i$, and a pure action selection $a_i(\cdot)$. By
condition~\textnormal{(a)}, every kernel
$P_{i,C}^{\sigma_{-i},a_i,\mathrm{aug}}$ has strictly positive internal
reachability between any two internal states in at most $|\Xi_C|-1$ steps. Hence
the displayed quantity in \eqref{eq:kappa-eta} is strictly positive for each fixed
$\sigma_{-i}$.

The map
$\sigma_{-i}\mapsto P_{i,C}^{\sigma_{-i},a_i,\mathrm{aug}}$ is continuous by
\Cref{lem:coeff-affine}, and the compact domain
$\prod_{j\ne i}\Sigma_{j,C}$ is a finite product of simplices. Therefore the
minimum over that domain is attained and remains strictly positive. Finally there
are only finitely many nodes, players, and pure action selections, so the overall
minimum is strictly positive.
\end{proof}

\subsection{Bias bound}


\begin{lemma}[Hitting the reference state or absorption]
\label{lem:hitting-time}
Let $X$ be a finite set, let $\partial$ be a finite absorbing set, and let $P$ be a
Markov kernel on $X\sqcup\partial$ such that for every $x,y\in X$ there exists
$1\le n\le |X|-1$ with $P^n(y\mid x)>0$. Define
\[
\kappa(P):=\min_{x,y\in X}\max_{1\le n\le |X|-1}P^n(y\mid x).
\]
Then for every $x,y\in X$,
\[
\E_x[\tau_{A_y}]\le \frac{|X|-1}{\kappa(P)},
\qquad
A_y:=\{y\}\cup \partial,
\qquad
\tau_{A_y}:=\inf\{t\ge 0:X_t\in A_y\}.
\]
\end{lemma}

\begin{proof}
Set $m:=|X|-1$. For each pair $(x,y)$ choose
$n(x,y)\in\{1,\dots,m\}$ such that
\[
P^{n(x,y)}(y\mid x)\ge \kappa(P).
\]
This is possible by the definition of $\kappa(P)$. Since $n(x,y)\le m$,
\[
\Pb_x(\tau_{A_y}\le m)\ge \Pb_x(X_{n(x,y)}=y)\ge \kappa(P).
\]
Indeed the event $\{X_{n(x,y)}=y\}$ already excludes prior absorption because all
states in $\partial$ are absorbing. Therefore
\[
\Pb_x(\tau_{A_y}>m)\le 1-\kappa(P).
\]
Applying the Markov property at times $m,2m,3m,\dots$ on the event
$\{\tau_{A_y}>km\}$, where the chain is still in $X$, gives
\[
\Pb_x(\tau_{A_y}>(k+1)m)\le (1-\kappa(P))\Pb_x(\tau_{A_y}>km)
\qquad (k\ge 0),
\]
hence
\[
\Pb_x(\tau_{A_y}>km)\le (1-\kappa(P))^k.
\]
Therefore
\[
\E_x[\tau_{A_y}]
=
\sum_{t\ge 0}\Pb_x(\tau_{A_y}>t)
\le
m\sum_{k\ge 0}\Pb_x(\tau_{A_y}>km)
\le
m\sum_{k\ge 0}(1-\kappa(P))^k
=
\frac{m}{\kappa(P)}.
\]
\end{proof}


\begin{proposition}[Uniform bias bound]
\label{prop:bias-span}
Assume condition~\textnormal{(a)}. Let $x^*\in\BR_\eta(x^*)$ be a fixed point and
let $(g_{i,C}^*,h_{i,C}^*)$ be the Bellman certificate of player $i$ in node $C$,
normalized by $h_{i,C}^*(\xi_C^\circ)=0$. Then
\[
\|h_{i,C}^*\|_\infty\le
H_\eta:=D_\eta\left(1+\frac{\max_D|\Xi_D|-1}{\kappa_\eta}\right).
\]
In particular,
\[
\spn(h_{i,C}^*)\le 2H_\eta.
\]
\end{proposition}

\begin{proof}
Fix $C$ and $i$. Let $a_i^*(\cdot)$ be a measurable selector choosing, in each
state, an action where equality holds in the Bellman system. Consider the genuine
augmented kernel
$P_{i,C}^{\sigma_{-i,C}^*,a_i^*,\mathrm{aug}}$ on
$\Xi_C^{\mathrm{aug}}$. Let
\[
A_C^\circ:=\{\xi_C^\circ\}\cup\partial B_C,
\qquad
\tau^\circ:=\inf\{t\ge 0:X_t\in A_C^\circ\}.
\]
By \Cref{lem:kappa-positive,lem:hitting-time},
\[
\E_\xi[\tau^\circ]\le \frac{|\Xi_C|-1}{\kappa_\eta}
\qquad (\xi\in\Xi_C).
\]

Along the selected policy $a_i^*(\cdot)$, the Bellman equalities read
\[
g_{i,C}^*+h_{i,C}^*(\xi)
=
r_{i,C}^{\sigma_{-i,C}^*}(\xi,a_i^*(\xi))
+
\sum_{\xi'}Q_{i,C}^{\sigma_{-i,C}^*}(\xi'\mid\xi,a_i^*(\xi))h_{i,C}^*(\xi')
+
\sum_b\Lambda_{i,C}^{\sigma_{-i,C}^*}(b\mid\xi,a_i^*(\xi))v_{i,C}^{x^*}(b).
\]
Stopping this identity at $\tau^\circ$ and using $h_{i,C}^*(\xi_C^\circ)=0$ and
the cemetery convention $h_{i,C}^*(\partial_b)=0$ gives the standard relative-value
representation
\[
h_{i,C}^*(\xi)
=
\E_\xi\!\left[
\sum_{t=0}^{\tau^\circ-1}\bigl(u_i^t-g_{i,C}^*\bigr)
+
\sum_{b\in B_C}\1_{\{X_{\tau^\circ}=\partial_b\}}\,v_{i,C}^{x^*}(b)
\right].
\]
Since $u_i^t\in[0,1]$, $g_{i,C}^*\in[0,d_\eta(C)]$, and
$v_{i,C}^{x^*}(b)\in[0,d_\eta(C)-1]$, the absolute value of the last display is
bounded by
\[
\begin{aligned}
d_\eta(C)\,\E_\xi[\tau^\circ] + d_\eta(C)-1
&\le
d_\eta(C)\bigl(\E_\xi[\tau^\circ]+1\bigr)
\\
&\le
D_\eta\left(1+\frac{|\Xi_C|-1}{\kappa_\eta}\right)
\\
&\le
D_\eta\left(1+\frac{\max_D|\Xi_D|-1}{\kappa_\eta}\right)
=
H_\eta.
\end{aligned}
\]
This proves the sup-norm bound. The span bound is immediate.
\end{proof}

\subsection{Deviation-value Lipschitz bound}

\begin{lemma}[Coefficient-level Bellman transfer]
\label{lem:coefficient-transfer}
Let $r,Q$ and $r',Q'$ be two finite average-reward MDP coefficient systems on the
same state-action graph, and let $(g,h)$ satisfy
\[
g+h(x)\ge r(x,a)+\sum_y Q(y\mid x,a)h(y)
\qquad\text{for all }(x,a).
\]
Then
\[
g+h(x)
\ge
r'(x,a)+\sum_y Q'(y\mid x,a)h(y)
-\Bigl(\|r'-r\|_\infty + \spn(h)\sup_{x,a}\|Q'(\cdot\mid x,a)-Q(\cdot\mid x,a)\|_{\mathrm{TV}}\Bigr)
\]
for all $(x,a)$. In particular the perturbation constant is governed by
$1+\spn(h)$, not merely by $\spn(h)$.
\end{lemma}

\begin{proof}
Fix $(x,a)$. Write
\begin{align*}
&r'(x,a)+\sum_y Q'(y\mid x,a)h(y)
-
\Bigl(r(x,a)+\sum_y Q(y\mid x,a)h(y)\Bigr)
\\
&= \bigl(r'(x,a)-r(x,a)\bigr)
+
\sum_y\bigl(Q'(y\mid x,a)-Q(y\mid x,a)\bigr)h(y).
\end{align*}
The first term is bounded by $\|r'-r\|_\infty$. For the second term, subtract the
minimum of $h$ and use that signed measures of total mass $0$ have dual norm equal
to total variation against functions of oscillation $\spn(h)$:
\[
\left|\sum_y\bigl(Q'-Q\bigr)(y\mid x,a)h(y)\right|
\le
\spn(h)\,\|Q'(\cdot\mid x,a)-Q(\cdot\mid x,a)\|_{\mathrm{TV}}.
\]
Combining the two bounds gives the claim.
\end{proof}

\begin{lemma}[Deviation-value Lipschitz bound]
\label{lem:deviation-lipschitz}
Fix $\eta\in(0,1)$ and a node $C$. There exists a finite constant $L_{\eta,C}$
such that for every player $i$, every two opponents' kernels
$\sigma_{-i},\tau_{-i}\in\prod_{j\ne i}\Sigma_{j,C}$, and every continuation vector
$v\in[0,d_\eta(C)-1]^{B_C\times N}$,
\[
\bigl|\operatorname{Val}_{i,C}(\tau_{-i},v)-\operatorname{Val}_{i,C}(\sigma_{-i},v)\bigr|
\le L_{\eta,C}\,\|\tau_{-i}-\sigma_{-i}\|_1,
\]
where $\operatorname{Val}_{i,C}(\cdot,v)$ denotes the optimal average reward of the
deviation MDP. One may take
\[
L_{\eta,C}=K_{i,C}(1+2H_\eta),
\]
with $K_{i,C}$ as in \Cref{lem:coeff-affine}.
\end{lemma}

\begin{proof}
Fix $\sigma_{-i}$ and let $(g,h)$ be an optimal Bellman certificate for the MDP with
coefficients induced by $(\sigma_{-i},v)$. By \Cref{prop:bias-span},
$\spn(h)\le 2H_\eta$. Apply \Cref{lem:coefficient-transfer} with the primed
coefficients induced by $(\tau_{-i},v)$. Using the Lipschitz estimates from
\Cref{lem:coeff-affine} yields
\[
\operatorname{Val}_{i,C}(\tau_{-i},v)
\le
\operatorname{Val}_{i,C}(\sigma_{-i},v)
+
K_{i,C}(1+2H_\eta)\,\|\tau_{-i}-\sigma_{-i}\|_1.
\]
Interchanging the roles of $\sigma_{-i}$ and $\tau_{-i}$ gives the reverse
inequality.
\end{proof}


\subsection{Blockwise Bellman inequality}

Fix a fixed point $x^*\in\BR_\eta(x^*)$ and the realized profile
$\sigma_{\eta,\delta,M_*}$ from \Cref{thm:realization}. For a complete block $B$
played in node $C$, let $X_{\mathrm{in}}(B)$ denote its entry phase-state and let
$X_{\mathrm{out}}(B)\in\Xi_C^{\mathrm{aug}}$ denote its terminal augmented state.
Extend each bias $h_{i,C}^*$ to $\Xi_C^{\mathrm{aug}}$ by the convention
\[
h_{i,C}^*(\partial_b):=0
\qquad (b\in B_C).
\]

\begin{lemma}
\label{lem:block-bellman}
There exists a constant $C_\eta<\infty$ such that for every complete block $B$
with cap $L(B)\ge M_*$ played in node $C$, every player $i$, and every unilateral
deviation $\tau_i$,
\begin{equation}
\label{eq:block-bellman-ineq}
\begin{aligned}
\E\left[
\widetilde U_i(B;\tau_i,\sigma_{-i}) - g_{i,C}^*|B|
\Bigm|\cF_{\mathrm{in}(B)}
\right]
&\le
(\eta+C_\eta\delta)\,\E\bigl[|B|\mid \cF_{\mathrm{in}(B)}\bigr]
\\
&\quad+
h_{i,C}^*(X_{\mathrm{in}}(B))
-
\E\bigl[h_{i,C}^*(X_{\mathrm{out}}(B))\mid\cF_{\mathrm{in}(B)}\bigr].
\end{aligned}
\end{equation}
Here $\widetilde U_i(B;\tau_i,\sigma_{-i})$ denotes the augmented block payoff when
player $i$ deviates inside block $B$ and the others keep the realized profile.
\end{lemma}

\begin{proof}
Fix $B$, $C$, and $i$. The Bellman certificate at the fixed point gives, for every
phase-state $\xi$ and every action $a_i\in A_i(\xi)$,
\begin{align}
\label{eq:exact-bellman-fixed-point}
g_{i,C}^* + h_{i,C}^*(\xi)
\ge {}
&r_{i,C}^{\sigma_{-i,C}^*}(\xi,a_i)
+
\sum_{\xi'\in\Xi_C}Q_{i,C}^{\sigma_{-i,C}^*}(\xi'\mid\xi,a_i)h_{i,C}^*(\xi')
\\
&+
\sum_{b\in B_C}\Lambda_{i,C}^{\sigma_{-i,C}^*}(b\mid\xi,a_i)v_{i,C}^{x^*}(b)
-\eta.
\nonumber
\end{align}
Inside a complete tail block, the compliant players are run by the controller
$\mathcal R_{C,L(B)}^{\sigma_C^*}$. By \Cref{ass:controller-input}, at every live
stage with current phase-state $\xi$ the conditional mixed action of the
nondeviators is exactly $\sigma_{-i,C}^*(\cdot\mid\xi)$. Hence the deviator sees
the exact one-step coefficients
$r_{i,C}^{\sigma_{-i,C}^*}$, $Q_{i,C}^{\sigma_{-i,C}^*}$, and
$\Lambda_{i,C}^{\sigma_{-i,C}^*}$ from
\eqref{eq:exact-bellman-fixed-point}; there is no separate coefficient-perturbation
term in the Bellman summation.

Let $\tau(B)=|B|$ be the actual stopping time of the block. Define the augmented
one-step reward $\widetilde u_i^t(\tau_i,\sigma_{-i})$ along the live block so that
it equals the ordinary stage payoff inside $\Xi_C$ and includes the continuation
term exactly at the live stage where the block first enters a cemetery state.
Summing the Bellman inequality from the block entry up to time $\tau(B)$ and
taking conditional expectations gives
\begin{align*}
\E\left[\widetilde U_i(B;\tau_i,\sigma_{-i})\Bigm|\cF_{\mathrm{in}(B)}\right]
\le{}&
\bigl(g_{i,C}^*+\eta\bigr)\,
\E\bigl[|B|\mid\cF_{\mathrm{in}(B)}\bigr]
\\
&+
h_{i,C}^*(X_{\mathrm{in}}(B))
-
\E\bigl[h_{i,C}^*(X_{\mathrm{out}}(B))\mid\cF_{\mathrm{in}(B)}\bigr].
\end{align*}
This is stronger than \eqref{eq:block-bellman-ineq}, since $C_\eta\delta\ge 0$.
The terminal term is zero on a cemetery exit by the extension convention above,
and it is the carried-over bias of the terminal internal state when the block ends
at its cap.
\end{proof}

\begin{corollary}[Augmented block regret]
\label{cor:block-regret}
There exists a constant $C_\eta<\infty$ such that for every complete block $B$
with cap at least $M_*$ played in node $C$, every player $i$, and every unilateral
deviation $\tau_i$,
\begin{equation}
\label{eq:block-regret-ineq}
\begin{aligned}
\E\left[
\widetilde U_i(B;\tau_i,\sigma_{-i})
-
\widetilde U_i(B;\sigma_{\eta,\delta,M_*})
\Bigm|\cF_{\mathrm{in}(B)}
\right]
&\le
(\eta+C_\eta\delta)\,\E\bigl[|B|\mid\cF_{\mathrm{in}(B)}\bigr]
\\
&\quad+
h_{i,C}^*(X_{\mathrm{in}}(B))
-
\E\bigl[h_{i,C}^*(X_{\mathrm{out}}(B))\mid\cF_{\mathrm{in}(B)}\bigr].
\end{aligned}
\end{equation}
\end{corollary}

\begin{proof}
Apply \Cref{lem:block-bellman} to the deviator's augmented block payoff. By
\eqref{eq:block-payoff-realization},
\[
\E\bigl[\widetilde U_i(B;\sigma_{\eta,\delta,M_*})\mid\cF_{\mathrm{in}(B)}\bigr]
\ge
\bigl(g_{i,C}^*-C_\eta\delta\bigr)\,
\E\bigl[|B|\mid\cF_{\mathrm{in}(B)}\bigr]
\]
after enlarging $C_\eta$ if necessary. Subtracting this lower bound from
\eqref{eq:block-bellman-ineq} yields \eqref{eq:block-regret-ineq}.
\end{proof}

\subsection{Proof of the main theorem}



\begin{proof}[Proof of \Cref{thm:main}]
Fix $\eps>0$. Choose parameters in the order
\[
\eta,\qquad \delta,\qquad M_*,\qquad T_0.
\]

\smallskip
\noindent\textbf{Step 1: fixed point.}
Choose $\eta>0$ with $\eta\le \eps/3$. By
\Cref{thm:kakutani-global}, there exists a fixed point $x^*\in\BR_\eta(x^*)$.

\smallskip
\noindent\textbf{Step 2: bias bound.}
Recall from \Cref{rem:tree-depth} that $D_\eta<\infty$. By
\Cref{lem:kappa-positive}, $\kappa_\eta>0$. Define
\[
H_\eta:=D_\eta\left(1+\frac{\max_D|\Xi_D|-1}{\kappa_\eta}\right).
\]
By \Cref{prop:bias-span}, every fixed-point Bellman bias satisfies
$\|h_{i,C}^*\|_\infty\le H_\eta$.

\smallskip
\noindent\textbf{Step 3: realization accuracy.}
Choose $\delta>0$ such that $C_\eta\delta\le \eps/3$, where $C_\eta$ is the
constant from \Cref{cor:block-regret}. Since the tree is finite, there is a common
cap threshold
\[
L_\eta(\delta):=\max_{C\in\cT_\eta}L_C(\delta)
\]
from \Cref{thm:local-controller}. Freeze a superexponential schedule at a level
$R$ so large that the associated tail cap $M_*=M_R$ satisfies
$M_*\ge L_\eta(\delta)$. Let
$\sigma^\eps:=\sigma_{\eta,\delta,M_*}$ be the realized profile from
\Cref{thm:realization}.

\smallskip
\noindent\textbf{Step 4: horizon threshold.}
By \Cref{lem:aggregation}, the telescoping remainder over complete tail blocks is
bounded by $2(D_\eta+1)H_\eta$, and the augmented/ordinary payoff discrepancy over
complete blocks is bounded by $2D_\eta^2$. The incomplete initial and terminal
fragments contribute at most a finite constant depending only on the frozen
prelude, the tail cap, and the finite tree object. Absorb all of these terms into
a constant $K_\eta<\infty$. Choose $T_0$ so large that
\[
\frac{K_\eta}{T_0}\le \frac{\eps}{3}.
\]

Now fix $T\ge T_0$, an initial state $s$, a player $i$, and a unilateral deviation
$\tau_i$. Let $B_1,\dots,B_m$ be the complete tail blocks contained in $[1,T]$,
listed in chronological order. The finitely many prelude blocks (before the
superexponential schedule reaches cap $M_*$) contribute at most $K_\eta$ to the
total payoff difference, hence at most $K_\eta/T$ to the average. This is
absorbed into the $K_\eta/T$ remainder term. Apply \Cref{cor:block-regret} block
by block to the
tail blocks and sum. Within each maximal run of consecutive blocks that stay in
the same node, the bias terms telescope because the terminal internal state of one
capped block is exactly the entry state of the next block in that same node. Every
time the node changes, the play moves to a strict child, and this can happen at
most $D_\eta$ times. Therefore \Cref{lem:aggregation} gives
\[
\frac1T\E\left[\sum_{r=1}^m
\bigl(
\widetilde U_i(B_r;\tau_i,\sigma^\eps_{-i})
-
\widetilde U_i(B_r;\sigma^\eps)
\bigr)\right]
\le
\eta + C_\eta\delta + \frac{K_\eta}{T}.
\]

It remains to pass from augmented block payoffs to ordinary stage payoffs. A block
differs from its ordinary stage-payoff sum only if it exits to a cemetery state,
in which case one continuation term in $[0,D_\eta]$ is added. Under either the
deviating or the compliant profile, the path can contain at most $D_\eta$ such
exits. Hence the total difference between augmented regret and ordinary regret over
all complete blocks in $[1,T]$ is at most $2D_\eta^2$, already absorbed into
$K_\eta$. The incomplete fragments are also already absorbed into $K_\eta$. We
obtain
\[
\gamma_i^T\bigl(s,(\tau_i,\sigma^\eps_{-i})\bigr)
-
\gamma_i^T(s,\sigma^\eps)
\le
\eta + C_\eta\delta + \frac{K_\eta}{T}.
\]
By construction, the right-hand side is at most
\[
\frac{\eps}{3}+\frac{\eps}{3}+\frac{\eps}{3}=\eps.
\]
Therefore
\[
\gamma_i^T(s,\sigma^\eps)
\ge
\gamma_i^T\bigl(s,(\tau_i,\sigma^\eps_{-i})\bigr)-\eps.
\]
Since $T\ge T_0$, $s$, $i$, and $\tau_i$ were arbitrary, $\sigma^\eps$ is a
public finite-memory uniform $\eps$-equilibrium.
\end{proof}


\section{Discussion}

The full Neyman conjecture remains open. What the present paper does is narrower
but mathematically sharp: it shows that the downstream compact-tree route is
correct once the finite universal $\eta$-tree object and its first-exit controller
realization are available and once the nodewise irreducibility obstruction from
condition~\textnormal{(a)} is excluded.

The central conceptual distinction in v9 is the separation between bookkeeping and
execution. The restart kernel $\rho_C$ belongs only to the convex bookkeeping
package; it is not an executed transition after a cemetery hit. Execution takes
place on the genuine absorbing augmented kernel, blocks terminate at the first
cemetery state, and the relevant local identity is renewal-reward: the bookkeeping
rates determine $g_{i,C}^*$, while the executed controller contributes that gain
through total augmented reward divided by actual first-exit time. With that
distinction in place, the fixed-point domain is convex, the Bellman target matches
the realized first-exit block, and the routing semantics are well-defined.

\appendix

\section{Finite average-reward MDP facts}
\label{app:mdp}

This appendix records the one-player linear-program facts used in the main text.
All statements are standard; see Puterman's monograph \cite{Puterman1994} for
background.

Consider a finite average-reward MDP with state space $X$, action sets $A(x)$,
bounded reward function $r(x,a)\in[0,\overline r]$ for some $\overline r<\infty$,
and transition kernel $Q(\cdot\mid x,a)$.

\begin{proposition}[Primal-dual formulation]
\label{prop:mdp-duality}
The optimal average reward is the common value of the primal and dual linear
programs
\begin{align*}
\textnormal{maximize }&\sum_{x\in X}\sum_{a\in A(x)}\nu(x,a)r(x,a)
\\
\textnormal{subject to }&\nu\ge 0,
\qquad \sum_{x,a}\nu(x,a)=1,
\\
&\sum_{a\in A(y)}\nu(y,a)=\sum_{x,a}\nu(x,a)Q(y\mid x,a)
\qquad (y\in X),
\end{align*}
and
\begin{align*}
\textnormal{minimize }&g
\\
\textnormal{subject to }&g+h(x)\ge r(x,a)+\sum_y Q(y\mid x,a)h(y)
\qquad (x\in X,\ a\in A(x)).
\end{align*}
There exists an optimal stationary mixed policy and an optimal dual pair $(g,h)$.
Moreover, after the linear normalization $h(x^\circ)=0$ at any fixed reference
state $x^\circ\in X$, one may choose $h$ in a compact box.
\end{proposition}

\begin{proof}
This is the classical average-reward LP formulation for finite MDPs. Feasibility is
immediate. The primal feasible set is a nonempty compact polytope, so an optimum
exists. The dual feasible set is nonempty because choosing $g=\overline r$ and
$h\equiv 0$ is feasible. Finite-dimensional LP duality gives equality of primal
and dual
optima and existence of optimal primal and dual solutions. The occupation measure
of an optimal primal solution induces an optimal stationary mixed policy. The
normalization $h(x^\circ)=0$ is obtained by subtracting the constant $h(x^\circ)$.
\end{proof}

\begin{proposition}[Uniqueness of the normalized bias in the communicating case]
\label{prop:bias-unique}
If the finite average-reward MDP is communicating, then any two optimal bias
vectors differ by a constant. Consequently, after fixing one reference coordinate
$h(x^\circ)=0$, the optimal bias is unique.
\end{proposition}

\begin{proof}
This is the standard communicating-MDP relative-value theorem; see
\cite[Chapter~8]{Puterman1994}. In the communicating case the optimal average
reward is state-independent and every optimal bias vector is unique up to an
additive constant. Fixing one coordinate removes that one-dimensional freedom.
\end{proof}


\section{Aggregating block estimates}
\label{app:aggregation}

\begin{lemma}
\label{lem:aggregation}
Let $B_1,\dots,B_m$ be the complete blocks contained in $[1,T]$, listed in
chronological order. For each $r$, let $C_r$ be the node of $B_r$, let
$\tau_r:=|B_r|$ be its actual duration, and let $\Delta_r$ be a block contribution.
Assume there exist constants $a\ge 0$, $H\ge 0$, $D\in\N$, and $K<\infty$ such that:
\begin{enumerate}[label=\textnormal{(\roman*)},leftmargin=2.6em]
\item $\sum_{r=1}^m \tau_r\le T$;
\item for every $r$ there exist $\cF_{\mathrm{in}(B_r)}$-measurable random variables
$\Psi_r$ and $\widetilde\Psi_r$ with $|\Psi_r|,|\widetilde\Psi_r|\le H$ a.s. and
\[
\E\bigl[\Delta_r\mid \cF_{\mathrm{in}(B_r)}\bigr]
\le
a\,\E\bigl[\tau_r\mid \cF_{\mathrm{in}(B_r)}\bigr]
+
\Psi_r-\E\bigl[\widetilde\Psi_r\mid \cF_{\mathrm{in}(B_r)}\bigr];
\]
\item whenever $B_r$ and $B_{r+1}$ lie in the same node, one has
$\widetilde\Psi_r=\Psi_{r+1}$ almost surely;
\item the number of indices $r$ for which $C_{r+1}\neq C_r$ is at most $D$;
\item the incomplete fragments outside the complete blocks contribute at most $K$
in expectation.
\end{enumerate}
Then
\[
\frac1T\E\left[\sum_{r=1}^m \Delta_r\right]
\le a + \frac{2(D+1)H+K}{T}.
\]
\end{lemma}

\begin{proof}
Partition $B_1,\dots,B_m$ into maximal consecutive runs on which the node is
constant. By assumption~(iv), there are at most $D+1$ such runs. Summing the
conditional estimates from assumption~(ii) and then taking expectations gives
\[
\E\left[\sum_{r=1}^m \Delta_r\right]
\le
a\sum_{r=1}^m \E[\tau_r]
+
\sum_{\text{runs}}
\bigl(\text{initial potential} - \text{terminal potential}\bigr).
\]
Within each run, the intermediate potentials cancel by assumption~(iii). Each run
therefore leaves at most one initial and one terminal term, each bounded in
absolute value by $H$. Because there are at most $D+1$ runs,
\[
\E\left[\sum_{r=1}^m \Delta_r\right]
\le
a\sum_{r=1}^m \E[\tau_r] + 2(D+1)H.
\]
Using assumption~(i) and adding the incomplete-fragment bound from assumption~(v)
yields
\[
\E\left[\sum_{r=1}^m \Delta_r\right]
\le
aT + 2(D+1)H + K.
\]
Divide by $T$.
\end{proof}

\begin{thebibliography}{99}

\bibitem{BewleyKohlberg1976}
T.~Bewley and E.~Kohlberg.
\newblock The asymptotic theory of stochastic games.
\newblock \emph{Mathematics of Operations Research}, 1(3):197--208, 1976.

\bibitem{FleschThuijsmanVrieze1997}
J.~Flesch, F.~Thuijsman, and K.~Vrieze.
\newblock Cyclic {M}arkov equilibria in stochastic games.
\newblock \emph{International Journal of Game Theory}, 26:303--314, 1997.

\bibitem{Kakutani1941}
S.~Kakutani.
\newblock A generalization of {B}rouwer's fixed point theorem.
\newblock \emph{Duke Mathematical Journal}, 8(3):457--459, 1941.

\bibitem{MertensNeyman1981}
J.-F.~Mertens and A.~Neyman.
\newblock Stochastic games.
\newblock \emph{International Journal of Game Theory}, 10(2):53--66, 1981.

\bibitem{MertensSorinZamir2003}
J.-F.~Mertens, S.~Sorin, and S.~Zamir.
\newblock Stochastic games.
\newblock In R.~Aumann and S.~Hart, editors, \emph{Handbook of Game Theory with
Economic Applications}, volume~3, chapter~47, pages 1811--1832. Elsevier,
Amsterdam, 2003.

\bibitem{Neyman2003}
A.~Neyman.
\newblock From {M}arkov chains to stochastic games.
\newblock In A.~Neyman and S.~Sorin, editors, \emph{Stochastic Games and
Applications}, NATO Science Series C, volume~570, pages 9--25. Kluwer Academic
Publishers, Dordrecht, 2003.

\bibitem{Nowak2003}
A.~S.~Nowak.
\newblock N-person stochastic games: Extensions of the finite state space case and
correlation.
\newblock In A.~Neyman and S.~Sorin, editors, \emph{Stochastic Games and
Applications}, NATO Science Series C, volume~570, pages 93--106. Kluwer Academic
Publishers, Dordrecht, 2003.

\bibitem{Puterman1994}
M.~L.~Puterman.
\newblock \emph{Markov Decision Processes: Discrete Stochastic Dynamic
Programming}.
\newblock John Wiley \& Sons, New York, 1994.

\bibitem{Shapley1953}
L.~S.~Shapley.
\newblock Stochastic games.
\newblock \emph{Proceedings of the National Academy of Sciences of the United
States of America}, 39(10):1095--1100, 1953.

\bibitem{Solan1999}
E.~Solan.
\newblock Three-player absorbing games.
\newblock \emph{Mathematics of Operations Research}, 24(3):669--698, 1999.

\bibitem{Solan2003}
E.~Solan.
\newblock Uniform equilibrium: More than two players.
\newblock In A.~Neyman and S.~Sorin, editors, \emph{Stochastic Games and
Applications}, NATO Science Series C, volume~570, pages 285--294. Kluwer Academic
Publishers, Dordrecht, 2003.

\bibitem{SolanVieille2002}
E.~Solan and N.~Vieille.
\newblock Correlated equilibrium in stochastic games.
\newblock \emph{Games and Economic Behavior}, 38(2):362--399, 2002.

\bibitem{Vieille2000I}
N.~Vieille.
\newblock Equilibrium in two-person stochastic games {I}: A reduction.
\newblock \emph{Israel Journal of Mathematics}, 119:55--91, 2000.

\bibitem{Vieille2000II}
N.~Vieille.
\newblock Equilibrium in two-person stochastic games {II}: The case of recursive
games.
\newblock \emph{Israel Journal of Mathematics}, 119:93--126, 2000.

\bibitem{Vieille2002}
N.~Vieille.
\newblock Stochastic games: Recent results.
\newblock In R.~Aumann and S.~Hart, editors, \emph{Handbook of Game Theory with
Economic Applications}, volume~3, chapter~48, pages 1833--1850. Elsevier,
Amsterdam, 2002.

\end{thebibliography}

\end{document}
